\documentclass[12pt]{article}
\usepackage[utf8]{inputenc}
\usepackage[T1]{fontenc}
\usepackage{hyperref}
\title{Karawari Language Resources for K-12 Curriculum Development}
\author{Compiled by Claude Sonnet 5 (claude-sonnet-5) for Andr\'as Kornai}
\date{}
\begin{document}
\maketitle

\section{Introduction}
This report covers Karawari (ISO 639-3 tzx, Glottocode tabr1243), a Lower Sepik language of East Sepik Province, Papua New Guinea. It is also called Tabriak or Yokoim. Speaker counts in the sources vary widely: 1,200 (Endangered Languages Project), 1,530 (Ethnologue 2015) and over 3,000 (Telban's fieldwork). Glottolog classes it as threatened. The Wamut language, spoken in the same river region, is a separate language.

\section{Primary Educational Resources}
I found no primers, textbooks or school programmes for Karawari in the search results.

\section{Linguistic and Reference Materials}
\begin{itemize}
\item Borut Telban, ``Saying, seeing, and knowing among the Karawari of Papua New Guinea'', a chapter in Aikhenvald and Dixon's \emph{The Grammar of Knowledge}. The record is listed as restricted access. \url{https://researchonline.jcu.edu.au/38488/}
\item Glottolog entry for Tabriak (Karawari). \url{https://glottolog.org/resource/languoid/id/tabr1243}
\item Wikipedia entry for the Tabriak language. \url{https://en.wikipedia.org/wiki/Tabriak_language}
\end{itemize}

\section{Religious and Community Resources}
I found none in the search results.

\section{Digital and Audio Resources}
\begin{itemize}
\item The Living Tongues Institute East Sepik project page describes the Yokoim (Karawari) variety, spoken in Kundiman, Manjamai and Konmei, and is linked to a Tabriak Talking Dictionary. I did not see the dictionary page itself. \url{https://livingtongues.org/east-sepik/}
\end{itemize}

\section{Curriculum Development Context}
The language is not known to be taught in schools, and I found no curriculum materials. The Talking Dictionary and the Living Tongues work are the most practical starting points. Any curriculum would have to be built with the community.

\end{document}